\documentclass[11pt,a4paper]{article}
\usepackage[utf8]{inputenc}
\usepackage[T1]{fontenc}
\usepackage{lmodern}
\usepackage[spanish,es-noquoting]{babel}
\usepackage[a4paper,margin=2.5cm]{geometry}
\usepackage{enumitem}
\usepackage{booktabs,array}
\usepackage[colorlinks=false,hidelinks]{hyperref}
\setlist{itemsep=1pt,topsep=2pt,parsep=0pt,leftmargin=1.2em}
\setlength{\parindent}{0pt}
\setlength{\parskip}{3pt}

% --- Datos del postulante (editar aquí) ---
\newcommand{\autor}{Máximo Sangiácomo}
\newcommand{\trayectoria}{Ingresé al BCRA en febrero de 2006: fui analista sénior y jefe de Investigaciones Económicas (crédito y finanzas corporativas) y jefe de Monitoreo Macroprudencial (2017--2020: regulación bancaria, crédito y estabilidad financiera), con personal a cargo en ambas jefaturas; desde mayo de 2020 soy Subgerente de Modelos Macroeconómicos: modelos DSGE y semiestructurales, \emph{pipelines} de pronóstico (series de tiempo y aprendizaje automático) y flujos de datos. Soy licenciado en Economía (UNLP) y magíster en Finanzas (UTDT), con especialización en Ciencia de Datos (ITBA), doctorado en Economía en curso (UNLP) e inglés avanzado.}

\title{\textbf{Proyecto de trabajo}\\[2pt]\large Gerente de Análisis Monetario (426 05)}
\author{\autor}
\date{}

\begin{document}
\maketitle

\section*{1. Introducción y enfoque}
Propongo que la Gerencia entregue tres resultados: pronósticos de liquidez y de efectivo mínimo evaluados ex post para la Subgerencia General de Operaciones, un diagnóstico monetario y crediticio oportuno y reproducible, y un equipo con procesos documentados y cobertura cruzada.

\textbf{Trayectoria.} \trayectoria{} En cada eje indico qué aporto y cómo cubro lo que aún no he hecho.

\textbf{Principios.} \emph{Continuidad antes que cambio}: ningún proceso diario se reemplaza sin correr antes en paralelo. \emph{Medir el error}: todo pronóstico se evalúa ex post. \emph{Robustez al régimen operativo}: las herramientas sirven cualquiera sea la combinación vigente de agregados, tasas y tipo de cambio. \emph{Adaptar, no copiar}: prácticas de otros bancos centrales, ajustadas al BCRA.

\textbf{Prioridades y responsables.} El \textbf{núcleo de los primeros 12 meses} son cuatro frentes: (1) el pronóstico de liquidez, efectivo mínimo y billetes, a cargo de Liquidez (Eje 1); (2) la evolución del informe monetario con agregados, FCI, billeteras y ECC, a cargo de Coyuntura (Ejes 2 y 3); (3) la programación monetaria trimestral y los módulos del IPOM, que coordino (Eje 4); y (4) la agenda de instrumentos, compartida (Eje 4). Código e IA (Eje 5) y la conducción del equipo (Eje 6) son habilitantes. Los frentes \emph{condicionados}, listados en la hoja de ruta, se activan cuando el ensamble de liquidez lleve tres meses en producción y haya horas disponibles; durante la corrida en paralelo no roto personal de Liquidez. Los pedidos urgentes ingresan por mí a un tablero de prioridades y desplazan lo condicionado, nunca la tarea diaria.

\textbf{Punto de partida.} Según su comunicado de política monetaria del 15 de diciembre de 2025, el BCRA enmarca su política en un esquema de agregados monetarios con banda de flotación cambiaria; publica además un informe monetario mensual, el IPOM trimestral y la Encuesta de Condiciones Crediticias (ECC). El diagnóstico de 90 días fija líneas de base y ajusta prioridades.

\section*{2. Ejes de trabajo}

\subsection*{Eje 1. Liquidez del mercado monetario y efectivo mínimo (funciones 1 y 5)}
\textbf{Qué}: un sistema de pronóstico de liquidez y de posición de efectivo mínimo con horizontes diario, semanal y mensual, por factores autónomos (sector público, sector externo, billetes), entregado a Operaciones en formato y horario acordados.
\textbf{Cómo}: ensamble de modelos para la posición de efectivo mínimo (exigencia e integración por entidad y período de cómputo) y para la liquidez (regresión con efectos de calendario---pagos, aguinaldo, vencimientos---, espacio de estado y ARIMAX) contra un benchmark ingenuo, con intervalos, y un modelo de circulación de billetes. De Hamilton (1996) tomo el peso del calendario en la demanda diaria de reservas, adaptado al período de cómputo del efectivo mínimo.
\textbf{Cómo se mide}: error absoluto medio relativo al benchmark, sesgo por día del mes y cobertura de los intervalos; tres meses de funcionamiento en paralelo antes de pasar a producción. Si a los nueve meses no corre en paralelo, sigue el proceso vigente, se extiende el plazo, reevalúo el umbral con Operaciones y postergo lo condicionado.
\textbf{Tasas de corto plazo}: Liquidez agrega al pronóstico un cuadro diario de call interbancario, caución, tasa de referencia y letras cortas del Tesoro, con sus diferenciales frente al excedente de liquidez proyectado; va a Operaciones y alimenta el \emph{briefing} y el informe.
\textbf{Aporte y brecha}: traigo el diseño de \emph{pipelines} de pronóstico. No he operado el pronóstico diario: en 90 días sigo con Liquidez y Operaciones el ciclo completo; durante los primeros seis meses las decisiones técnicas de liquidez llevan firma conjunta con su subgerente, y ningún cambio en producción se decide sin su aval y el de sus jefes.

\subsection*{Eje 2. Agregados monetarios y coyuntura (funciones 2 y 6)}
El BCRA ya publica un informe monetario mensual; el objetivo es \emph{evolucionarlo}, no crearlo. \emph{Núcleo}: descomposición de la base monetaria y del M2 por componentes y fuentes de creación, con desestacionalización y ajuste por calendario con software abierto (X-13/JDemetra+), generada desde el código. \emph{Condicionado}: nota quincenal de dos páginas; agregados Divisia (Barnett, 1980), como los índices de servicios monetarios que publicó la Fed de St.~Louis; y bases con \emph{vintages}, al estilo de ALFRED.
\textbf{Aporte y brecha}: no he elaborado el informe de coyuntura monetaria. En 90 días acompaño dos ciclos completos del informe con el/la Subgerente de Coyuntura y su equipo, armamos una ficha por producto (destinatario, uso, insumos, horas) y acordamos qué se mantiene, simplifica o cambia; luego, revisión trimestral.
\textbf{Presentaciones y discursos (función 6)}: desde un banco de gráficos generado por el código, Coyuntura entrega láminas y mensajes clave en 72 horas (24 si es urgente), con cifras trazables y revisión del jefe, del subgerente y mía antes de elevarlas.

\subsection*{Eje 3. Crédito, medios de pago y mercado de capitales (función 3 y perfil)}
\textbf{Crédito}: la ECC se integra (\emph{núcleo}) al informe y al IPOM (índices de difusión, poder predictivo sobre el crédito), como el \emph{Bank Lending Survey} del BCE, sin crear otra encuesta. \emph{Condicionado}: panel de microdatos de la Central de Deudores inspirado en el Consumer Credit Panel de la Fed de Nueva York (flujos, transiciones de situación), \emph{sujeto a un acuerdo de acceso y resguardo con Estadísticas y Superintendencia}. \textbf{Medios de pago y mercado de capitales}: la pregunta es cuánta liquidez migra de depósitos a caución, FCI de mercado de dinero y billeteras, y cómo eso altera los agregados y el pronóstico de liquidez. \emph{Núcleo}: seguimiento mensual de FCI de mercado de dinero, caución, transferencias y QR, con datos públicos y del área que los produce. \emph{Condicionado}: \emph{perímetro del dinero}, cuadro mensual que estima qué parte de esos activos y de los títulos cortos funciona como reserva de liquidez frente a M2. No cubro valuación de activos ni análisis de emisores. \textbf{Aporte}: catorce años de análisis del crédito en el BCRA (2006--2020) y trabajos sobre relaciones bancarias (\emph{IMF Economic Review}, 2020) y exportadores (BIS \emph{Working Papers}, 2020).

\subsection*{Eje 4. Programación monetaria e instrumentos (funciones 4, 5 y 7)}
\textbf{Programación (función 4)}: un ejercicio trimestral en tres bloques. El primero estima la demanda de dinero, para la base y el M2 privado, en función de la actividad, la inflación esperada y el costo de oportunidad, con un modelo de corrección de errores en ventanas móviles que se compara con uno simple de velocidad. El segundo asegura la consistencia: proyecta las fuentes de la base según el balance del BCRA, incluido el resultado cuasifiscal, las concilia con el programa financiero del Tesoro y el balance cambiario, y agrega préstamos y depósitos privados por segmento, moneda y grupo de entidades, consistentes con los encajes; los escenarios llevan bandas de incertidumbre en formato \emph{fan chart}, como en el Bank of England (BoE). El tercero sigue los desvíos: una banda de tolerancia para base y M2, derivada del error histórico de la demanda, se revisa cada mes, y si el desvío la excede elevo una nota que separa demanda de fuentes y propone reestimar, ajustar instrumentos o revisar el programa.
\textbf{Instrumentos (función 5)}: además de responder consultas, propongo una agenda propia. Cada año, una nota técnica evalúa ex post las medidas de encajes e instrumentos de liquidez con estudios de eventos sobre caución y curva corta y comparando entidades según su posición de efectivo mínimo previa. Cada semestre se estima el traspaso de la tasa de referencia y de la caución a las tasas pasivas y activas por línea. Con esos resultados se calibran las notas de respuesta a consultas, que exponen canal de transmisión, efecto sobre liquidez, agregados y costo cuasifiscal, y riesgos: en horas si son urgentes, en días si son estándar.
\textbf{IPOM (función 7)}: módulos y recuadros con responsables y cronograma acordado con la Gerencia Principal (borrador a $T-15$ días, cierre a $T-5$). \textbf{Aporte y brecha}: los escenarios usan, por acuerdo con el área de modelos, los modelos DSGE y semiestructurales que conozco; la programación operativa parte de los ejercicios vigentes.

\subsection*{Eje 5. Datos, código abierto e inteligencia artificial (perfil y función 10)}
\textbf{Infraestructura}: repositorio versionado en entorno interno (Python y R), informes generados desde el código, pruebas de datos y revisión por pares. \textbf{Aprendizaje automático}: regularización y \emph{boosting} como complemento, contra benchmarks lineales y con técnicas de interpretabilidad (SHAP, importancia por permutación); el BoE ya usó aprendizaje automático en pronóstico (Chakraborty y Joseph, 2017). \textbf{IA generativa}: el primer caso de uso, en entorno interno, es la asistencia en el borrador del informe mensual y en la revisión de código; se mide por las horas de producción del informe y los defectos detectados en revisión, con la meta de bajar las horas sin aumentar los defectos. \textbf{Gobernanza}: la política de IA (mes 6) define qué datos pueden usarse (públicos y agregados sí; datos por entidad, de liquidez o de política nunca salen del entorno interno), registra cada uso, exige revisión humana de todo producto y valida con Sistemas y Seguridad de la Información todo modelo en producción. \textbf{Aporte}: programo en Stata, R y Python, publiqué once módulos de Stata y dicté el curso interno de R del BCRA (2019 y 2021).

\subsection*{Eje 6. Conducción, equipo y articulación institucional (funciones 8 a 11)}
\begin{itemize}
\item \textbf{Planificación y delegación (funciones 9 y 11)}: plan anual con objetivos medibles por subgerencia (Análisis de Coyuntura Monetaria; Análisis y Gestión de la Liquidez) y por persona, seguido cada trimestre; diagnóstico de dotación y competencias en 90 días; cada subgerente tiene presupuesto de horas y autoridad sobre sus frentes, reunión individual mensual y evaluación de desempeño; sucesión y doble cobertura; reuniones semanales de coyuntura y mensuales de calidad de pronósticos.
\item \textbf{Desarrollo de colaboradores}: mentoría de analistas y asistentes (enseño en la UNLP desde 2010), capacitación individual (CCBS del BoE, donde cursé, FMI, CEMLA) con transferencia al equipo, reconocimiento, rotación y publicación.
\item \textbf{Articulación}: \emph{briefing} semanal de una página (liquidez, agregados, crédito) elevado por la Gerencia Principal; acuerdo de servicio con Operaciones y calendario con Estadísticas.
\item \textbf{Confidencialidad y ética}: la información de liquidez y de política mueve mercados; por eso, acceso según necesidad, perfiles segregados, registro de accesos revisado cada semestre con Seguridad de la Información, restricción de difusión antes de cada decisión, salida sólo por canales formales y declaración de conflictos de interés.
\item \textbf{Auditoría (función 8)}: tablero de observaciones con responsable, plazo y evidencia, revisado cada trimestre hasta que la Auditoría General da cada una por cerrada.
\end{itemize}

\section*{3. Hoja de ruta e indicadores}
\small
\begin{tabular}{@{}p{2.4cm}p{12.4cm}@{}}
\toprule
\textbf{0--3 meses} & Diagnóstico (ciclo de liquidez, dos ciclos del informe, fichas, línea de base de errores, capacidad); plan anual; acuerdo con Operaciones; repositorio.\\
\textbf{3--12 meses} & \emph{Núcleo}: (1) ensamble de liquidez y billetes en paralelo y luego en producción; (2) informe con agregados desestacionalizados, FCI, billeteras y ECC; (3) primera programación con banda de desvíos; (4) protocolo de consultas, traspaso y evaluación ex post (mes 12). Habilitante: política de IA.\\
\textbf{12--24 meses} & Informe anual de calidad de pronósticos y segunda evaluación ex post; \emph{condicionado}: nota quincenal, Divisia y \emph{vintages} (Eje 2), panel de crédito y perímetro del dinero (Eje 3); dos documentos de trabajo; revisión de la gestión.\\
\bottomrule
\end{tabular}

\normalsize\textbf{Capacidad (supuesto)}: Liquidez, operación diaria 55\,\% y paralelo 25\,\%; Coyuntura, informe e IPOM 50\,\% y evolución e instrumentos 30\,\%; ambas, programación 10\,\% y reserva 10\,\%.

\textbf{Indicadores y metas}: (i) a los 12 meses, error del ensamble de liquidez menor que el del benchmark según el test de Diebold y Mariano (1995) al 10\,\% con al menos 60 observaciones diarias en paralelo; si no se reúnen, reducción del error con intervalo que excluya el cero; (ii) procesos críticos con manual y doble cobertura: 50\,\% a 6 meses, 100\,\% a 9; (iii) al menos 95\,\% de los informes y módulos del IPOM entregados en fecha; (iv) en 24 meses, dos documentos de trabajo y cuatro notas técnicas (dos, evaluaciones ex post); (v) 40 horas anuales de capacitación por persona; (vi) cada desvío trimestral de base o M2 fuera de la banda, explicado en nota dentro de diez días hábiles, separando error de pronóstico de decisiones de política; (vii) valoración anual de Operaciones y autoridades, mínimo 4 sobre 5.

\normalsize
\section*{4. Riesgos y mitigación}
\small
\begin{tabular}{@{}>{\raggedright\arraybackslash}p{3.0cm}>{\raggedright\arraybackslash}p{2.9cm}>{\raggedright\arraybackslash}p{6.4cm}>{\raggedright\arraybackslash}p{2.3cm}@{}}
\toprule
\textbf{Riesgo} & \textbf{Señal} & \textbf{Mitigación} & \textbf{Responsable}\\
\midrule
Cambio de régimen o de autoridades & Quiebre en errores & Reestimación, escenarios alternativos, plan revisado en 30 días & Gerente\\
Falla o ausencia en liquidez & Error fuera del intervalo & Doble cobertura, manuales, modelo previo; aviso inmediato y causas en 48 h & Subg.\ de Liquidez\\
Datos y Sistemas & Demoras en insumos & Controles automáticos; calendarios con Estadísticas y Sistemas & Subg.\ de Coyuntura\\
Sobrecarga & Horas sobre lo presupuestado & Se posterga lo condicionado y luego el frente 4 (salvo consultas) & Ambos subgerentes\\
Resistencia & Baja participación & Diseño participativo y pilotos & Gerente\\
\bottomrule
\end{tabular}
\normalsize

\small
\section*{Referencias}
Barnett (1980), \emph{J.\ Econometrics} 14(1). Chakraborty y Joseph (2017), \emph{Machine learning at central banks}, BoE SWP 674. Diebold y Mariano (1995), \emph{JBES} 13(3). Hamilton (1996), \emph{JPE} 104(1).
\end{document}
